\documentclass[12pt]{article}
\usepackage[utf8]{inputenc}
\usepackage[T1]{fontenc}
\usepackage{hyperref}
\title{Tzotzil Language Resources for K-12 Curriculum Development}
\author{Compiled by Claude Sonnet 5 (claude-sonnet-5) for Andr\'as Kornai}
\date{}
\begin{document}
\maketitle
\section{Introduction}
Tzotzil (Bats'i k'op, tzo) is a Mayan language of the Chiapas highlands in Mexico, with more than 400,000 speakers. This report is based on one round of searches, so it is not exhaustive.
\section{Primary Educational Resources}
The SEP (Mexican Ministry of Education) published \emph{Bats'i k'op} reading and exercise books for indigenous primary schools. Library catalogs list a 1993 first-cycle reader: \url{https://opac.biblio.iteso.mx/vufind/Record/000499546/Channels} and \url{https://sibi.upn.mx/bib/201580}. Some catalog entries list the series under Tseltal rather than Tsotsil, so check the items. A scanned third- and fourth-grade reader is online in the SEP historical collection: \url{https://dgeiib.basica.sep.gob.mx/files/fondo-editorial/acervo-historico/ah_00141.pdf}.
\section{Linguistic and Reference Materials}
Robert Laughlin's \emph{Great Tzotzil Dictionary of San Lorenzo Zinacant\'an} (Smithsonian, 1975) has about 30,000 entries and is open access: \url{https://repository.si.edu/items/34078bcc-1288-422c-9c1a-50b81954a440}. The multilingual dictionary \emph{Svunal bats'i k'opetik} covers Spanish, Tseltal, Tsotsil, Ch'ol and Tojolobal: \url{https://baul.fahho.mx/Record/27562/Holdings}. A Living Dictionary is at \url{https://livingdictionaries.app/batsi-kop-tsotsil-tsot/synopsis}.
\section{Religious and Community Resources}
A PDF of Tzotzil language material from a Bible project is at \url{https://www.iglesiacatolica.org.gt/bibliakiche/BKiche_Lenguas_LenguaTsotsil.pdf}. I did not verify its contents beyond the background text.
\section{Digital and Audio Resources}
Apart from the online dictionaries above, none were found.
\section{Curriculum Development Context}
Bilingual materials exist from the SEP primary series, but they date from the 1990s and are not known to be used widely. Secondary-level materials were not found. The Laughlin dictionary covers the Zinacant\'an variety, and other varieties differ.
\end{document}